\documentclass[12pt, twoside]{article}
\input{../preamble}

\fancyhead[LO]{La Scuola d'Italia / Dr.\ Huson / IB Math \\* 29 May 2026}
\fancyhead[RO]{Markscheme}
\fancyfoot[C]{\thepage}

\begin{document}
\subsubsection*{8.6 Early Finishers: Transformations of Quadratics and a Cubic --- Solutions}

\medskip

%================================================================
%--- Problem 1
%================================================================
\noindent\textbf{1.} $\displaystyle g(x)=-2(x+3)^2+5$, \ parent $f(x)=x^2$.

\medskip

\textbf{(a)} From $y=x^2$, in order:

\quad reflection in the $x$-axis (the factor $-2$ is negative);

\quad vertical stretch with scale factor $2$;

\quad translation $3$ units left and $5$ units up
$\left(\text{vector } \left(\begin{smallmatrix}-3\\5\end{smallmatrix}\right)\right)$.

\emph{Accept the reflection and stretch described together as ``vertical
stretch with scale factor $-2$,'' and in either order relative to each other.}

\medskip

\textbf{(b)} Vertex $(-3,\,5)$; \ axis of symmetry $x=-3$.

\medskip

\textbf{(c)} The parabola opens downward, so the vertex is a maximum:
range $y\le 5$ \ (i.e. $\{\,y\in\mathbb{R} : y\le 5\,\}$).

\medskip

\textbf{(d)} $g(x)=-2(x^2+6x+9)+5=-2x^2-12x-18+5=-2x^2-12x-13$.

\medskip

\textbf{(e)} Downward parabola, vertex $(-3,5)$, $y$-intercept
$g(0)=-2(9)+5=-13$, so $(0,-13)$.

$x$-intercepts (for reference): $(x+3)^2=\tfrac{5}{2}$, so
$x=-3\pm\sqrt{\tfrac{5}{2}}\approx -1.42,\ -4.58$.

\bigskip\hrule\bigskip

%================================================================
%--- Problem 2
%================================================================
\noindent\textbf{2.} $y=x^2$: \ vertical stretch $\times\tfrac{1}{2}$
$\rightarrow$ reflection in $x$-axis $\rightarrow$ translation
$\left(\begin{smallmatrix}4\\-3\end{smallmatrix}\right)$.

\medskip

\textbf{(a)} Stretch: $y=\tfrac{1}{2}x^2$. \ Reflect in $x$-axis:
$y=-\tfrac{1}{2}x^2$. \ Translate $4$ right and $3$ down:

$\displaystyle h(x)=-\tfrac{1}{2}(x-4)^2-3$.

\medskip

\textbf{(b)} Vertex $(4,\,-3)$.

No $x$-intercepts: the parabola opens downward and its maximum value is
$-3<0$, so the graph lies entirely below the $x$-axis. (Equivalently,
$-\tfrac{1}{2}(x-4)^2-3=0$ gives $(x-4)^2=-6$, which has no real solution.)

\medskip

\textbf{(c)} No --- the order does not matter. Both a vertical stretch
($\times\tfrac{1}{2}$) and a reflection in the $x$-axis ($\times(-1)$) are
vertical scalings, and they multiply the output: $(-1)\cdot\tfrac{1}{2}
=\tfrac{1}{2}\cdot(-1)=-\tfrac{1}{2}$. Multiplication is commutative, so either
order yields $y=-\tfrac{1}{2}x^2$.

\bigskip\hrule\bigskip

%================================================================
%--- Problem 3
%================================================================
\noindent\textbf{3.} $\displaystyle g(x)=2(x+1)^3-16$, \ parent $f(x)=x^3$.

\medskip

\textbf{(a)} From $y=x^3$, in order:

\quad vertical stretch with scale factor $2$;

\quad translation $1$ unit left and $16$ units down
$\left(\text{vector } \left(\begin{smallmatrix}-1\\-16\end{smallmatrix}\right)\right)$.

\medskip

\textbf{(b)} The point of inflection is the image of the origin:
$(0,0)\mapsto(-1,\,-16)$.

\medskip

\textbf{(c)} $x$-intercept: $2(x+1)^3-16=0$, \ $(x+1)^3=8$, \ $x+1=2$, \ $x=1$,
so $(1,0)$.

$y$-intercept: $g(0)=2(1)^3-16=-14$, so $(0,-14)$.

\medskip

\textbf{(d)} Increasing cubic with point of inflection $(-1,-16)$, passing
through $(1,0)$ and $(0,-14)$.

\bigskip\hrule\bigskip

%================================================================
%--- Problem 4 (Challenge)
%================================================================
\noindent\textbf{4.} \textbf{Challenge.}

\medskip

\textbf{(a)} Complete the square:

$p(x)=-x^2+8x-13=-(x^2-8x)-13=-\big[(x-4)^2-16\big]-13=-(x-4)^2+3$.

From $y=x^2$: a reflection in the $x$-axis, then a translation $4$ units right
and $3$ units up $\left(\text{vector } \left(\begin{smallmatrix}4\\3\end{smallmatrix}\right)\right)$.

Vertex $(4,\,3)$ (a maximum).

\medskip

\textbf{(b)} The point of inflection of $y=a(x-h)^3+k$ is at $(h,k)$, so
$h=-1$ and $k=2$.

Substitute $(1,18)$: \ $18=a(1-(-1))^3+2=a(2)^3+2=8a+2$.

$8a=16$, \ $a=2$.

Equation: $\displaystyle y=2(x+1)^3+2$.

\end{document}
